\chapter{Introducción}\label{cap:introduccion}

\section{Propósito del documento}

Fijar \emph{qué} tiene que hacer el proyecto antes de seguir construyéndolo: los
requerimientos de los datos, de la aplicación web y de calidad, y la estructura con
la que se organiza el trabajo. Es la referencia común del equipo: si algo no está
aquí, se discute y se añade; si cambia, se cambia aquí primero.

Documentos relacionados:

\begin{tabularx}{\linewidth}{@{}lX@{}}
\toprule
Documento & Contenido \\
\midrule
\texttt{README.md} & Resumen del reto y del problema con nuestros datos. \\
\texttt{analisis/serie-armonizada.md} & Método de armonización y su validación. \\
\texttt{Documentacion/avance.tex} & Estado del proyecto a la fecha. \\
\texttt{Documentacion/frontend/plan.tex} & Propuesta inicial del frontend. \\
\bottomrule
\end{tabularx}

\section{Contexto: el reto}

Reto \emph{Harmonization of MODIS and VIIRS Hot Spots} del NASA Space Apps Challenge
2026, propuesto por la CONAE (Comisión Nacional de Actividades Espaciales, Argentina).

Los focos de calor que publica NASA FIRMS vienen de dos familias de sensores: MODIS
(desde 2000) y VIIRS (desde 2012). VIIRS detecta muchos más focos que MODIS sobre el
mismo territorio, así que una serie que empalma los dos muestra un aumento de incendios
que no existe. El reto pide resolver esa incompatibilidad y ofrecer el resultado de
forma útil.

\section{Objetivo}

Construir un \textbf{calendario de actividad de quemas} consistente desde el año 2000
hasta hoy, en el que el número de focos signifique lo mismo sin importar el satélite que
los detectó, y ofrecerlo en una \textbf{aplicación web} con la que una persona no
experta pueda responder en menos de un minuto:

\begin{itemize}
  \item ¿Este mes es peor de lo normal en mi zona?
  \item ¿Cuál fue el peor año, y en qué meses se concentra la actividad?
  \item ¿Cuántos focos hubo en esta área y periodo? (y descargarlos)
\end{itemize}

\section{Usuarios}

\begin{tabularx}{\linewidth}{@{}lX@{}}
\toprule
Usuario & Qué necesita \\
\midrule
Gestores de emergencias y bomberos & Saber rápido si el periodo actual es anómalo. \\
Gestores del territorio & Comparar zonas y épocas; planificar temporadas. \\
Científicos & Series consistentes, incertidumbre explícita y descarga de datos. \\
Jurado del reto & Entender el problema y la solución en pocos minutos. \\
\bottomrule
\end{tabularx}

\section{Alcance}

\textbf{Dentro del MVP:}
\begin{itemize}
  \item Armonización MODIS $\to$ VIIRS para 2000--2025 con validación documentada.
  \item Aplicación web con mapa, selección de área y periodo, calendario
        años~$\times$~meses, comparación con lo normal, rankings y descarga en CSV.
  \item Una región de trabajo (India, porque ya tenemos sus datos).
\end{itemize}

\textbf{Fuera del MVP:}
\begin{itemize}
  \item Predicción de incendios.
  \item Datos casi en tiempo real (NRT). Deseable, pero secundario.
  \item Cuentas de usuario.
\end{itemize}

\abierta{\textbf{cobertura mundial.} El README la deja fuera del MVP, pero se ha
planteado mostrar datos de todo el mundo. Implica descargar el archivo global de FIRMS
($\sim$7\,GB comprimidos), recalcular la constante de conversión por región (la actual
solo vale para India) y servir los datos de otra forma (sección~\ref{sec:decisiones}).
Los requerimientos se escriben de modo que la región sea configurable
(RF-01), para que esta decisión no obligue a rehacer la aplicación.}

\section{Convenciones}

Cada requerimiento tiene un identificador y una prioridad:

\begin{tabularx}{\linewidth}{@{}lX@{}}
\toprule
Prefijo & Tipo \\
\midrule
\texttt{RD} & Requerimiento de datos (capítulo~\ref{cap:datos}). \\
\texttt{RF} & Requerimiento funcional de la aplicación (capítulo~\ref{cap:funcionales}). \\
\texttt{RNF} & Requerimiento no funcional o de calidad (capítulo~\ref{cap:nofuncionales}). \\
\bottomrule
\end{tabularx}

\begin{tabularx}{\linewidth}{@{}lX@{}}
\toprule
Prioridad & Significado \\
\midrule
\imp & Sin esto no hay entrega. \\
\rel & Debe estar en la entrega; se recorta solo si no hay tiempo. \\
\des & Mejora si sobra tiempo. \\
\bottomrule
\end{tabularx}

\section{Glosario}

\begin{description}[style=nextline]
  \item[Foco de calor] Píxel en el que un satélite detecta una anomalía térmica.
        Normalmente un fuego activo, pero también puede ser una industria o un volcán.
  \item[FIRMS] \emph{Fire Information for Resource Management System}, el servicio de
        NASA que distribuye los focos de MODIS y VIIRS.
  \item[MODIS] Sensor de los satélites Terra (desde noviembre de 2000) y Aqua (desde
        julio de 2002). Píxel de $\sim$1\,km.
  \item[VIIRS] Sensor de S-NPP (2012), NOAA-20 (2018) y NOAA-21 (2023). Píxel de
        $\sim$375\,m; detecta fuegos más pequeños.
  \item[Armonizar] Llevar los focos de los dos sensores a una escala común. Aquí, la de
        VIIRS S-NPP.
  \item[Constante $k$] Factor por el que se multiplica cada foco MODIS para expresarlo
        en focos VIIRS equivalentes. Se calcula por mes del año.
  \item[Peso] Valor de cada foco en la escala armonizada: $k$ para MODIS, 1 para VIIRS.
        Los focos armonizados de un área son la suma de los pesos.
  \item[Lo normal] Valor de referencia de un mes o año para un área (mediana o promedio
        histórico; decisión abierta).
  \item[Anomalía] Diferencia entre un periodo y lo normal, en porcentaje.
\end{description}
